% =============================================================================
%  CHAPTER 5: Mathematics
% =============================================================================

\chapter{Mathematics}

\latex{} is renowned for its mathematical typesetting. This chapter covers everything from basic inline formulas to advanced constructs.

% ------------------------------------------------------------------
\section{Math Modes}

There are several ways to enter math mode:

\begin{table}[H]
\centering
\begin{tabular}{lll}
\toprule
\textbf{Type} & \textbf{Syntax} & \textbf{Usage} \\
\midrule
Inline math   & \verb|$...$|            & $E=mc^2$ in text \\
Display math  & \verb|\[...\]|          & Centered, numbered \\
Equation env  & \verb|\begin{equation}| & Centered, numbered \\
Align env     & \verb|\begin{align}|    & Multi-line aligned \\
\bottomrule
\end{tabular}
\caption{Math mode entry points}
\end{table}

\begin{lstlisting}[language={[LaTeX]TeX}, caption=Math modes]
% Inline: the formula flows with text
The energy is $E = mc^2$ according to Einstein.

% Display: the formula is centered on its own line
\[ E = mc^2 \]

% Equation: like display but with a number
\begin{equation}
    E = mc^2 \label{eq:einstein}
\end{equation}

% Starred version: no number
\begin{equation*}
    E = mc^2
\end{equation*}
\end{lstlisting}

% ------------------------------------------------------------------
\section{Superscripts and Subscripts}

\begin{lstlisting}[language={[LaTeX]TeX}, caption=Superscripts and subscripts]
$x^2$              % Superscript: x squared
$x_i$              % Subscript: x sub i
$x_i^2$            % Both: x sub i squared
$x_{i,j}$          % Multi-character subscript: need braces
$a^{b^c}$          % Nested superscript
$2^{10} = 1024$    % Powers
\end{lstlisting}

Results: $x^2$, \quad $x_i$, \quad $x_i^2$, \quad $x_{i,j}$, \quad $a^{b^c}$, \quad $2^{10} = 1024$

% ------------------------------------------------------------------
\section{Fractions}

\begin{lstlisting}[language={[LaTeX]TeX}, caption=Fractions]
$\frac{a}{b}$                          % Basic fraction
$\dfrac{a}{b}$                         % Display-style fraction
$\tfrac{a}{b}$                         % Text-style fraction
$\cfrac{1}{1+\cfrac{1}{2}}$           % Continued fraction
$\binom{n}{k}$                         % Binomial coefficient
\end{lstlisting}

Results: $\frac{a}{b}$, \quad $\dfrac{a}{b}$, \quad $\cfrac{1}{1+\cfrac{1}{2}}$, \quad $\binom{n}{k}$

% ------------------------------------------------------------------
\section{Roots}

\begin{lstlisting}[language={[LaTeX]TeX}, caption=Roots]
$\sqrt{x}$            % Square root
$\sqrt[3]{x}$         % Cube root
$\sqrt[n]{x}$         % nth root
$\sqrt{a^2 + b^2}$    % Complex expression
\end{lstlisting}

Results: $\sqrt{x}$, \quad $\sqrt[3]{x}$, \quad $\sqrt[n]{x}$, \quad $\sqrt{a^2 + b^2}$

% ------------------------------------------------------------------
\section{Sums, Products, Integrals}

\begin{lstlisting}[language={[LaTeX]TeX}, caption=Big operators]
$\sum_{i=1}^{n} x_i$         % Summation
$\prod_{i=1}^{n} x_i$        % Product
$\int_{a}^{b} f(x) dx$       % Definite integral
$\iint_{D} f(x,y) dA$        % Double integral
$\iiint_{V} f(x,y,z) dV$     % Triple integral
$\oint_{C} F \cdot dr$       % Contour integral
$\lim_{x \to \infty} f(x)$  % Limit
$\sup_{x \in S} f(x)$       % Supremum
$\max_{x \in S} f(x)$       % Maximum
$\min_{x \in S} f(x)$       % Minimum
\end{lstlisting}

Display style results:
\[ \sum_{i=1}^{n} x_i \qquad \prod_{i=1}^{n} x_i \qquad \int_{a}^{b} f(x)\,dx \qquad \lim_{x \to \infty} f(x) \]

% ------------------------------------------------------------------
\section{Greek Letters}

\begin{table}[H]
\centering
\begin{tabular}{llll}
\toprule
\textbf{Command} & \textbf{Result} & \textbf{Command} & \textbf{Result} \\
\midrule
\verb|\alpha|     & $\alpha$     & \verb|\Alpha|     & $A$ \\
\verb|\beta|      & $\beta$      & \verb|\gamma|     & $\gamma$ \\
\verb|\delta|     & $\delta$     & \verb|\epsilon|   & $\epsilon$ \\
\verb|\varepsilon|& $\varepsilon$& \verb|\zeta|      & $\zeta$ \\
\verb|\eta|       & $\eta$       & \verb|\theta|     & $\theta$ \\
\verb|\iota|      & $\iota$      & \verb|\kappa|     & $\kappa$ \\
\verb|\lambda|    & $\lambda$    & \verb|\mu|        & $\mu$ \\
\verb|\nu|        & $\nu$        & \verb|\xi|        & $\xi$ \\
\verb|\pi|        & $\pi$        & \verb|\rho|       & $\rho$ \\
\verb|\sigma|     & $\sigma$     & \verb|\tau|       & $\tau$ \\
\verb|\upsilon|   & $\upsilon$   & \verb|\phi|       & $\phi$ \\
\verb|\varphi|    & $\varphi$    & \verb|\chi|       & $\chi$ \\
\verb|\psi|       & $\psi$       & \verb|\omega|     & $\omega$ \\
\midrule
\verb|\Gamma|     & $\Gamma$     & \verb|\Delta|     & $\Delta$ \\
\verb|\Theta|     & $\Theta$     & \verb|\Lambda|    & $\Lambda$ \\
\verb|\Xi|        & $\Xi$        & \verb|\Pi|        & $\Pi$ \\
\verb|\Sigma|     & $\Sigma$     & \verb|\Phi|       & $\Phi$ \\
\verb|\Psi|       & $\Psi$       & \verb|\Omega|     & $\Omega$ \\
\bottomrule
\end{tabular}
\caption{Greek letters}
\end{table}

% ------------------------------------------------------------------
\section{Mathematical Operators and Functions}

\begin{lstlisting}[language={[LaTeX]TeX}, caption=Math operators]
$\sin x$     $\cos x$     $\tan x$
$\arcsin x$  $\arccos x$  $\arctan x$
$\sinh x$    $\cosh x$    $\tanh x$
$\ln x$      $\log x$     $\exp x$
$\det A$     $\dim V$     $\ker f$
$\deg f$     $\gcd(a,b)$  $\hom(A,B)$
$\inf S$     $\sup S$     $\max S$
$\min S$     $\lim f$     $\limsup f$
\end{lstlisting}

% ------------------------------------------------------------------
\section{Relations and Binary Operators}

\begin{table}[H]
\centering
\begin{tabular}{llll}
\toprule
\textbf{Command} & \textbf{Symbol} & \textbf{Command} & \textbf{Symbol} \\
\midrule
\verb|<|      & $<$      & \verb|>|       & $>$ \\
\verb|\leq|   & $\leq$   & \verb|\geq|    & $\geq$ \\
\verb|\neq|   & $\neq$   & \verb|\approx| & $\approx$ \\
\verb|\equiv| & $\equiv$ & \verb|\sim|    & $\sim$ \\
\verb|\propto|& $\propto$& \verb|\ll|     & $\ll$ \\
\verb|\gg|    & $\gg$    & \verb|\in|     & $\in$ \\
\verb|\subset|& $\subset$& \verb|\supset| & $\supset$ \\
\verb|\subseteq| & $\subseteq$ & \verb|\supseteq| & $\supseteq$ \\
\verb|\cup|   & $\cup$   & \verb|\cap|    & $\cap$ \\
\verb|\setminus| & $\setminus$ & \verb|\emptyset| & $\emptyset$ \\
\verb|\pm|    & $\pm$    & \verb|\mp|     & $\mp$ \\
\verb|\times| & $\times$ & \verb|\div|    & $\div$ \\
\verb|\cdot|  & $\cdot$  & \verb|\circ|   & $\circ$ \\
\verb|\otimes|& $\otimes$& \verb|\oplus|  & $\oplus$ \\
\bottomrule
\end{tabular}
\caption{Common mathematical symbols}
\end{table}

% ------------------------------------------------------------------
\section{Arrows}

\begin{table}[H]
\centering
\begin{tabular}{llll}
\toprule
\textbf{Command} & \textbf{Arrow} & \textbf{Command} & \textbf{Arrow} \\
\midrule
\verb|\leftarrow|      & $\leftarrow$      & \verb|\rightarrow|     & $\rightarrow$ \\
\verb|\Leftarrow|      & $\Leftarrow$      & \verb|\Rightarrow|     & $\Rightarrow$ \\
\verb|\leftrightarrow| & $\leftrightarrow$ & \verb|\Leftrightarrow| & $\Leftrightarrow$ \\
\verb|\mapsto|         & $\mapsto$         & \verb|\longmapsto|     & $\longmapsto$ \\
\verb|\uparrow|        & $\uparrow$        & \verb|\downarrow|      & $\downarrow$ \\
\verb|\nearrow|        & $\nearrow$        & \verb|\searrow|        & $\searrow$ \\
\bottomrule
\end{tabular}
\caption{Arrow symbols}
\end{table}

% ------------------------------------------------------------------
\section{Matrices}

The \texttt{amsmath} package provides several matrix environments:

\begin{lstlisting}[language={[LaTeX]TeX}, caption=Matrices]
% Parentheses: ( )
\begin{pmatrix}
    a & b \\
    c & d
\end{pmatrix}

% Brackets: [ ]
\begin{bmatrix}
    a & b \\
    c & d
\end{bmatrix}

% Braces: { }
\begin{Bmatrix}
    a & b \\
    c & d
\end{Bmatrix}

% Vertical bars: | |
\begin{vmatrix}
    a & b \\
    c & d
\end{vmatrix}

% Double vertical bars: || ||
\begin{Vmatrix}
    a & b \\
    c & d
\end{Vmatrix}

% No delimiters
\begin{matrix}
    a & b \\
    c & d
\end{matrix}

% With dots:
\begin{pmatrix}
    a_{11} & a_{12} & \cdots & a_{1n} \\
    a_{21} & a_{22} & \cdots & a_{2n} \\
    \vdots & \vdots & \ddots & \vdots \\
    a_{m1} & a_{m2} & \cdots & a_{mn}
\end{pmatrix}
\end{lstlisting}

Results:
\[
\begin{pmatrix} a & b \\ c & d \end{pmatrix}
\quad
\begin{bmatrix} 1 & 2 \\ 3 & 4 \end{bmatrix}
\quad
\begin{vmatrix} a & b \\ c & d \end{vmatrix}
\]

\[
\begin{pmatrix}
a_{11} & a_{12} & \cdots & a_{1n} \\
a_{21} & a_{22} & \cdots & a_{2n} \\
\vdots & \vdots & \ddots & \vdots \\
a_{m1} & a_{m2} & \cdots & a_{mn}
\end{pmatrix}
\]

% ------------------------------------------------------------------
\section{Multi-line Equations (align)}

\begin{lstlisting}[language={[LaTeX]TeX}, caption=align environment]
\begin{align}
    f(x) &= x^2 + 2x + 1 \\
         &= (x+1)^2 \\
         &= x^2 + 2x + 1 \label{eq:expand}
\end{align}

% align without numbering:
\begin{align*}
    f(x) &= x^2 + 2x + 1 \\
         &= (x+1)^2
\end{align*}

% Number specific lines only:
\begin{align}
    f(x) &= x^2 + 2x + 1 \nonumber \\
         &= (x+1)^2 \\
         &= x^2 + 2x + 1 \nonumber
\end{align}

% Multiple alignment columns:
\begin{align*}
    x &= 1  &  y &= 2  &  z &= 3 \\
    x' &= 4 &  y' &= 5 &  z' &= 6
\end{align*}
\end{lstlisting}

Result:
\begin{align}
    f(x) &= x^2 + 2x + 1 \label{eq:quadratic} \\
         &= (x+1)^2 \nonumber
\end{align}

% ------------------------------------------------------------------
\section{Cases and Piecewise Functions}

\begin{lstlisting}[language={[LaTeX]TeX}, caption=cases environment]
\begin{equation}
    |x| =
    \begin{cases}
        x  & \text{if } x \geq 0 \\
        -x & \text{if } x < 0
    \end{cases}
\end{equation}
\end{lstlisting}

Result:
\begin{equation}
    |x| =
    \begin{cases}
        x  & \text{if } x \geq 0 \\
        -x & \text{if } x < 0
    \end{cases}
\end{equation}

% ------------------------------------------------------------------
\section{Accents in Math Mode}

\begin{table}[H]
\centering
\begin{tabular}{lll}
\toprule
\textbf{Command} & \textbf{Result} & \textbf{Name} \\
\midrule
\verb|\hat{a}|     & $\hat{a}$     & Hat \\
\verb|\bar{a}|     & $\bar{a}$     & Bar \\
\verb|\tilde{a}|   & $\tilde{a}$   & Tilde \\
\verb|\vec{a}|     & $\vec{a}$     & Vector arrow \\
\verb|\dot{a}|     & $\dot{a}$     & Dot \\
\verb|\ddot{a}|    & $\ddot{a}$    & Double dot \\
\verb|\vec{a}|     & $\vec{a}$     & Vector \\
\verb|\widehat{abc}|& $\widehat{abc}$& Wide hat \\
\verb|\widetilde{abc}|& $\widetilde{abc}$& Wide tilde \\
\verb|\overline{abc}| & $\overline{abc}$ & Overline \\
\verb|\underline{abc}| & $\underline{abc}$ & Underline \\
\verb|\overbrace{abc}| & $\overbrace{abc}$ & Overbrace \\
\verb|\underbrace{abc}|& $\underbrace{abc}$& Underbrace \\
\bottomrule
\end{tabular}
\caption{Math accents}
\end{table}

% ------------------------------------------------------------------
\section{Spacing in Math Mode}

\begin{table}[H]
\centering
\begin{tabular}{lll}
\toprule
\textbf{Command} & \textbf{Space} & \textbf{Example} \\
\midrule
\verb|\,|     & Thin space      & $a\,b$ \\
\verb|\:|     & Medium space    & $a\:b$ \\
\verb|\;|     & Thick space     & $a\;b$ \\
\verb|\!|     & Negative thin   & $a\!b$ \\
\verb|\quad|  & 1em space       & $a\quad b$ \\
\verb|\qquad| & 2em space       & $a\qquad b$ \\
\verb|\,|     & & $ab$ (no space for comparison) \\
\bottomrule
\end{tabular}
\caption{Math spacing commands}
\end{table}

% ------------------------------------------------------------------
\section{Text Inside Math Mode}

\begin{lstlisting}[language={[LaTeX]TeX}, caption=Text in math]
$\text{for all } x > 0$           % Regular text in math
$\text{if and only if}$            % Using amsmath \text
$\mathrm{d}x$                      % Upright d for differentials
$\text{Area} = \pi r^2$            % Upright labels
$\text{let } n = 5$                % Mixed text and math
\end{lstlisting}

% ------------------------------------------------------------------
\section{Theorems and Proofs}

\begin{lstlisting}[language={[LaTeX]TeX}, caption=Theorem environments]
% In preamble (amsthm package):
\newtheorem{theorem}{Theorem}[section]
\newtheorem{lemma}[theorem]{Lemma}
\newtheorem{corollary}[theorem]{Corollary}
\theoremstyle{definition}
\newtheorem{definition}{Definition}[section]
\theoremstyle{remark}
\newtheorem{remark}{Remark}

% In document:
\begin{theorem}[Pythagorean]
For a right triangle: $a^2 + b^2 = c^2$.
\end{theorem}

\begin{proof}
By construction of squares on each side...
\end{proof}
\end{lstlisting}

\newtheorem{theorem}{Theorem}[section]
\newtheorem{lemma}[theorem]{Lemma}
\theoremstyle{definition}
\newtheorem{definition}{Definition}[section]

\begin{theorem}[Pythagorean]
For a right triangle with legs $a$, $b$ and hypotenuse $c$:
\[ a^2 + b^2 = c^2 \]
\end{theorem}

\begin{proof}
Left as an exercise for the reader.
\end{proof}

\begin{definition}
A \textbf{prime number} is a natural number greater than 1 that has no positive divisors other than 1 and itself.
\end{definition}

% ------------------------------------------------------------------
\section{Comprehensive Math Example}

\begin{equation}
    \int_{-\infty}^{\infty} e^{-x^2}\,dx = \sqrt{\pi}
    \label{eq:gaussian}
\end{equation}

\begin{equation}
    e^{i\pi} + 1 = 0
    \label{eq:euler}
\end{equation}

The Gaussian integral (\ref{eq:gaussian}) and Euler's identity (\ref{eq:euler}) are two of the most beautiful equations in mathematics.
